% \iffalse meta-comment
%
% File: luacjk.dtx
% -----------------------------------------------------------------------
%   Copyright (C) 2026 by Mingyu Xia <myhsia@outlook.com>               *
% -----------------------------------------------------------------------
%   This work may be distributed and/or modified under the conditions   *
%   of the LaTeX Project Public License (LPPL), either version 1.3c of  *
%   this license or (at your option) any later version.                 *
%   The latest version of this license is in                            *
%                                                                       *
%       http://www.latex-project.org/lppl.txt                           *
%                                                                       *
%   and version 1.3c or later is part of all distributions of LaTeX     *
%   version 2008 or later.                                              *
%                                                                       *
%   This work has the LPPL maintenance status `maintained'.             *
%                                                                       *
%   The Current Maintainer of this work is Mingyu Xia.                  *
% -----------------------------------------------------------------------
%   This work consists of the files luacjk.dtx,                         *
%                                   luacjk.ins,                         *
%                 the derived files luacjk.sty,                         *
%                                   luacjk.lua,                         *
%           the documentation files luacjk.pdf,                         *
%                               and README.md.                          *
% -----------------------------------------------------------------------
%
%   Any modification of this file should ensure that the copyright and  *
%   license information is placed in the derived files.                 *
%                                                                       *
% -----------------------------------------------------------------------
%
% \fi
% \title{\bfseries The \pkg{LuaCJK} package}
% \author{Mingyu Xia\thanks{\url{https://github.com/myhsia/LuaCJK}}}
% \date{Released \DocFileDate\quad \texttt{\DocFileVersion}}
% \maketitle
%
% \begin{abstract}
%   \pkg{LuaCJK} provides Chinese typesetting for LuaLaTeX.  It is inspired
%   by \pkg{xeCJK}, but uses LuaTeX node callbacks instead of XeTeX
%   character classes.  This manual documents the first development release.
% \end{abstract}
%
% \tableofcontents
%
% \begin{documentation}
%
% \section{Introduction}
%
% \pkg{LuaCJK} is a LuaTeX-only package for Chinese typesetting.  Its public
% interface follows the familiar \pkg{xeCJK} commands:
% \cs{setCJKmainfont}, \cs{setCJKsansfont}, \cs{setCJKmonofont},
% \cs{setCJKfamilyfont}, \cs{CJKfamily}, and \cs{LuaCJKsetup}.
%
% The implementation does not use XeTeX's inter-character token mechanism.
% Instead, LuaCJK records the current CJK font in a LuaTeX attribute and
% replaces CJK glyph nodes before line breaking and box packing.
%
% \section{Loading}
%
% \begin{quote}
%   \cs{usepackage}\oarg{options}\marg{luacjk}
% \end{quote}
%
% LuaCJK requires LuaLaTeX.  It conflicts with \pkg{luatexja},
% \pkg{luatex-cn}, and \pkg{xeCJK}; loading one of these packages before or
% after LuaCJK raises an error.
%
% \section{Font selection}
%
% The main font commands take an optional list of \pkg{fontspec} options.
% Options may be written before or after the font name.
%
% \begin{function}{\setCJKmainfont,\setCJKsansfont,\setCJKmonofont}
%   \begin{syntax}
%     \tn{setCJKmainfont} \oarg{options} \marg{font} \oarg{options}
%   \end{syntax}
%   Set the CJK font for the roman, sans serif, and monospaced families.
% \end{function}
%
% \iffalse
%<*example>
% \fi
\begin{codemo}
{\CJKfamily{zhsong} 中文宋体} \quad
{\CJKfamily{zhhei} 中文黑体} \quad
{\CJKfamily{zhkai} 中文楷体} \quad
{\CJKfamily{zhfs} 中文仿宋}
\end{codemo}
% \iffalse
%</example>
% \fi
%
% \begin{function}{\setCJKfamilyfont}
%   \begin{syntax}
%     \tn{setCJKfamilyfont} \marg{family} \oarg{options} \marg{font} \oarg{options}
%   \end{syntax}
%   Define a named CJK family.
% \end{function}
%
% \begin{function}{\CJKfamily}
%   \begin{syntax}
%     \tn{CJKfamily} \marg{family}
%   \end{syntax}
%   Switch to a named CJK family.
% \end{function}
%
% \begin{function}{\newCJKfontfamily}
%   \begin{syntax}
%     \tn{newCJKfontfamily} \oarg{family} \meta{command} \oarg{options} \marg{font}
%   \end{syntax}
%   Define a command that switches to a newly declared family.  The command is
%   local to the current group, like \cs{newfontfamily}.
% \end{function}
%
% \begin{function}{\CJKfontspec}
%   \begin{syntax}
%     \tn{CJKfontspec} \oarg{options} \marg{font}
%   \end{syntax}
%   Select a CJK font for the current group.
% \end{function}
%
% \section{Spacing}
%
% \begin{function}{\LuaCJKsetup}
%   \begin{syntax}
%     \tn{LuaCJKsetup} \marg{key-value list}
%   \end{syntax}
%   Configure LuaCJK.  The settings are local to the current group.
% \end{function}
%
% \begin{function}{CJKglue}
%   \begin{syntax}
%     CJKglue = \meta{glue}
%   \end{syntax}
%   Glue inserted between adjacent CJK characters.  The initial value is
%   \texttt{0pt plus 0.08\cs{baselineskip}}.
% \end{function}
%
% \begin{function}{CJKecglue}
%   \begin{syntax}
%     CJKecglue = \meta{glue}
%   \end{syntax}
%   Glue inserted between CJK and Latin characters.  The initial value is a
%   normal interword space.
% \end{function}
%
% \begin{function}{CJKspace}
%   \begin{syntax}
%     CJKspace = \meta{true} \textbar{} \meta{false}
%   \end{syntax}
%   Whether to keep source spaces between CJK characters.  The initial value
%   is \meta{false}.
% \end{function}
%
% \begin{function}{xCJKecglue}
%   \begin{syntax}
%     xCJKecglue = \meta{true} \textbar{} \meta{false}
%   \end{syntax}
%   Whether to replace source spaces between CJK and Latin characters with
%   \texttt{CJKecglue}.  The initial value is \meta{false}.
% \end{function}
%
% \iffalse
%<*example>
% \fi
\begin{codemo}
中文 English 中文
\end{codemo}
% \iffalse
%</example>
% \fi
%
% \iffalse
%<*example>
% \fi
\begin{codemo}
\begingroup
  \LuaCJKsetup{CJKspace = true}
  中 文 中 文
\endgroup
\end{codemo}
% \iffalse
%</example>
% \fi
%
% \section{Current limitations}
%
% The first release intentionally covers only the font and spacing core.
% Punctuation compression, CJK mathematics, fallback fonts, alternate fonts,
% the decoration layer, and a \pkg{ctex} backend are not yet implemented.
%
% \end{documentation}
%
% \StopEventually{\clearpage\PrintIndex}
%
% \begin{implementation}
%
% \section{Implementation}
%
% \begin{macrocode}
%<*package>
%    \end{macrocode}
%
% \begin{macrocode}
%<@@=luacjk>
%    \end{macrocode}
%
% \begin{macrocode}
%<package>\GetIdInfo $Id: luacjk.dtx 0.1.0 2026-11-02 myhsia<myhsia@outlook.com>$
%<package>  {A LuaTeX CJK typesetting package}
%<package>\ProvidesExplPackage {\ExplFileName}
%<!driver>  {\ExplFileDate} {\ExplFileVersion} {\ExplFileDescription}
%<package>\edef \DocFileDate        {\ExplFileDate}
%<package>\edef \DocFileVersion     {\ExplFileVersion}
%<package>\edef \DocFileDescription {\ExplFileDescription}
%    \end{macrocode}
%
% \begin{macrocode}
\sys_if_engine_luatex:F
  {
    \msg_new:nnn { luacjk } { engine }
      { LuaCJK~requires~LuaTeX. }
    \msg_critical:nn { luacjk } { engine }
    \tex_endinput:D
  }
\RequirePackage { expl3 }
\RequirePackage { l3keys2e }
\RequirePackage { fontspec }
\RequirePackage { xparse }
%    \end{macrocode}
%
% \begin{macrocode}
\msg_new:nnn { luacjk } { luatexja-conflict }
  {
    LuaCJK~and~LuaTeX-ja~must~not~be~loaded~together.\\
    Load~only~one~of~them.
  }
\msg_new:nnn { luacjk } { luatexcn-conflict }
  {
    LuaCJK~and~LuaTeX-CN~must~not~be~loaded~together.\\
    Load~only~one~of~them.
  }
\msg_new:nnn { luacjk } { xecjk-conflict }
  {
    LuaCJK~is~for~LuaTeX;~xeCJK~is~for~XeTeX.\\
    Load~only~one~of~them.
  }
\cs_if_exist:cT { ver@luatexja.sty }
  { \msg_critical:nn { luacjk } { luatexja-conflict } }
\cs_if_exist:cT { ver@luatex-cn.sty }
  { \msg_critical:nn { luacjk } { luatexcn-conflict } }
\cs_if_exist:cT { ver@luatex-cn-hori.sty }
  { \msg_critical:nn { luacjk } { luatexcn-conflict } }
\cs_if_exist:cT { ver@xeCJK.sty }
  { \msg_critical:nn { luacjk } { xecjk-conflict } }
\AddToHook { package/luatexja/before }
  { \msg_critical:nn { luacjk } { luatexja-conflict } }
\AddToHook { package/luatex-cn/before }
  { \msg_critical:nn { luacjk } { luatexcn-conflict } }
\AddToHook { package/luatex-cn-hori/before }
  { \msg_critical:nn { luacjk } { luatexcn-conflict } }
\AddToHook { package/xeCJK/before }
  { \msg_critical:nn { luacjk } { xecjk-conflict } }
%    \end{macrocode}
%
% \begin{macrocode}
\newattribute \g_luacjk_font_attr
\int_const:Nn \c_luacjk_font_attr_int { \the \allocationnumber }
\tl_const:Nx \c_luacjk_version_tl { \ExplFileVersion }
\cs_new:Npn \LuaCJKversion { \tl_use:N \c_luacjk_version_tl }
%    \end{macrocode}
%
% \begin{macrocode}
\skip_new:N \l_luacjk_ccglue_skip
\skip_new:N \l_luacjk_ecglue_skip
\skip_set:Nn \l_luacjk_ccglue_skip { 0pt plus 0.08 \tex_baselineskip:D }
\skip_set:Nn \l_luacjk_ecglue_skip
  { 0.3333em plus 0.1667em minus 0.1111em }
\newattribute \g_luacjk_cjk_space_attr
\int_const:Nn \c_luacjk_cjk_space_attr_int { \the \allocationnumber }
\newattribute \g_luacjk_xecglue_attr
\int_const:Nn \c_luacjk_xecglue_attr_int { \the \allocationnumber }
\g_luacjk_cjk_space_attr = 0
\g_luacjk_xecglue_attr = 0
%    \end{macrocode}
%
% \begin{macrocode}
\keys_define:nn { luacjk / options }
  {
    CJKglue .code:n =
      {
        \skip_set:Nn \l_luacjk_ccglue_skip {#1}
      },
    CJKecglue .code:n =
      {
        \skip_set:Nn \l_luacjk_ecglue_skip {#1}
      },
    CJKspace .choice:,
    CJKspace / true  .code:n = { \g_luacjk_cjk_space_attr = 1 },
    CJKspace / false .code:n = { \g_luacjk_cjk_space_attr = 0 },
    CJKspace .default:n = { true },
    CJKspace .initial:n = { false },
    xCJKecglue .choice:,
    xCJKecglue / true  .code:n = { \g_luacjk_xecglue_attr = 1 },
    xCJKecglue / false .code:n = { \g_luacjk_xecglue_attr = 0 },
    xCJKecglue .default:n = { true },
    xCJKecglue .initial:n = { false },
  }
\NewDocumentCommand \LuaCJKsetup { m }
  { \keys_set:nn { luacjk / options } {#1} }
%    \end{macrocode}
%
% \begin{macrocode}
\prop_new:N \g_luacjk_family_prop
\tl_new:N \l_luacjk_family_tl
\tl_set:Nn \l_luacjk_family_tl { zhsong }
\tl_new:N \l_luacjk_tmpa_tl
\tl_new:N \l_luacjk_fontspec_tl
\tl_new:N \l_luacjk_family_name_tl
\tl_new:N \CJKrmdefault
\tl_new:N \CJKsfdefault
\tl_new:N \CJKttdefault
\tl_new:N \CJKfamilydefault
\tl_set:Nn \CJKrmdefault { zhsong }
\tl_set:Nn \CJKsfdefault { zhhei }
\tl_set:Nn \CJKttdefault { zhfs }
\tl_set:Nn \CJKfamilydefault { zhsong }
%    \end{macrocode}
%
% \begin{macrocode}
\cs_new_protected:Npn \luacjk_set_family:nnn #1#2#3
  {
    \group_begin:
      \fontspec_gset_family:Nnn \l_luacjk_fontspec_tl {#2} {#3}
      \prop_gput:NnV \g_luacjk_family_prop {#1} \l_luacjk_fontspec_tl
    \group_end:
    \prop_if_in:NVT \g_luacjk_family_prop \l_luacjk_family_tl
      { \selectfont }
  }
\cs_new_protected:Npn \luacjk_switch_family:n #1
  {
    \prop_if_in:NnTF \g_luacjk_family_prop {#1}
      {
        \tl_set:Nn \l_luacjk_family_tl {#1}
        \selectfont
      }
      { \msg_warning:nnn { luacjk } { family-unknown } {#1} }
  }
\msg_new:nnn { luacjk } { family-unknown }
  { Unknown~CJK~family~`#1'~is~being~ignored. }
%    \end{macrocode}
%
% \begin{macrocode}
\bool_new:N \l_luacjk_selectfont_bool
\int_new:N \g_luacjk_fontid_int
\cs_new_protected:Npn \luacjk_update_font:
  {
    \bool_if:NF \l_luacjk_selectfont_bool
      {
        \bool_set_true:N \l_luacjk_selectfont_bool
        \prop_get:NVNTF \g_luacjk_family_prop \l_luacjk_family_tl \l_luacjk_family_name_tl
          {
            \group_begin:
              \exp_args:NV \fontfamily \l_luacjk_family_name_tl
              \selectfont
              \int_gset:Nn \g_luacjk_fontid_int
                { \tex_fontid:D \tex_font:D }
            \group_end:
            \g_luacjk_font_attr = \g_luacjk_fontid_int
          }
          { \g_luacjk_font_attr = \c_zero_int }
        \bool_set_false:N \l_luacjk_selectfont_bool
      }
  }
\AddToHook { selectfont } { \luacjk_update_font: }
%    \end{macrocode}
%
% \begin{macrocode}
\NewDocumentCommand \setCJKmainfont { o m o }
  {
    \exp_args:NV \luacjk_set_family_with_options:nnnn
      \CJKrmdefault {#1} {#2} {#3}
  }
\NewDocumentCommand \setCJKsansfont { o m o }
  {
    \exp_args:NV \luacjk_set_family_with_options:nnnn
      \CJKsfdefault {#1} {#2} {#3}
  }
\NewDocumentCommand \setCJKmonofont { o m o }
  {
    \exp_args:NV \luacjk_set_family_with_options:nnnn
      \CJKttdefault {#1} {#2} {#3}
  }
\NewDocumentCommand \setCJKfamilyfont { m o m o }
  { \luacjk_set_family_with_options:nnnn {#1} {#2} {#3} {#4} }
\cs_new_protected:Npn \luacjk_set_family_with_options:nnnn #1#2#3#4
  {
    \tl_if_novalue:nTF {#2}
      {
        \tl_if_novalue:nTF {#4}
          { \luacjk_set_family:nnn {#1} { } {#3} }
          { \luacjk_set_family:nnn {#1} {#4} {#3} }
      }
      { \luacjk_set_family:nnn {#1} {#2} {#3} }
  }
\NewDocumentCommand \CJKfamily { m }
  { \luacjk_switch_family:n {#1} \tex_ignorespaces:D }
\NewDocumentCommand \newCJKfontfamily { o m o m }
  {
    \tl_set:Ne \l_luacjk_tmpa_tl
      {
        \IfNoValueTF {#1}
          { \cs_to_str:N #2 }
          {#1}
      }
    \cs_set_protected:Npx #2
      { \exp_not:N \luacjk_switch_family:n { \exp_not:V \l_luacjk_tmpa_tl } }
    \exp_args:NV \luacjk_set_family_with_options:nnnn
      \l_luacjk_tmpa_tl {#3} {#4} { \NoValue }
  }
\int_new:N \g_luacjk_fontspec_int
\NewDocumentCommand \CJKfontspec { o m }
  {
    \int_gincr:N \g_luacjk_fontspec_int
    \tl_set:Ne \l_luacjk_tmpa_tl
      { luacjk/fontspec/ \int_use:N \g_luacjk_fontspec_int }
    \exp_args:NV \luacjk_set_family_with_options:nnnn
      \l_luacjk_tmpa_tl {#1} {#2} { \NoValue }
    \exp_args:NV \luacjk_switch_family:n \l_luacjk_tmpa_tl
    \tex_ignorespaces:D
  }
%    \end{macrocode}
%
% \begin{macrocode}
\NewDocumentCommand \songti   { } { \CJKfamily { zhsong } }
\NewDocumentCommand \heiti    { } { \CJKfamily { zhhei  } }
\NewDocumentCommand \fangsong { } { \CJKfamily { zhfs   } }
\NewDocumentCommand \kaishu   { } { \CJKfamily { zhkai  } }
\NewDocumentCommand \lishu    { } { \CJKfamily { zhli   } }
\NewDocumentCommand \youyuan  { } { \CJKfamily { zhyou  } }
\AddToHook { cmd/rmfamily/after }
  { \exp_args:NV \luacjk_switch_family:n \CJKrmdefault }
\AddToHook { cmd/sffamily/after }
  { \exp_args:NV \luacjk_switch_family:n \CJKsfdefault }
\AddToHook { cmd/ttfamily/after }
  { \exp_args:NV \luacjk_switch_family:n \CJKttdefault }
%    \end{macrocode}
%
% \begin{macrocode}
\group_begin:
\char_set_catcode_space:n { 32 }
\directlua{
  local f = kpse.find_file("luacjk.lua")
  if f then
    dofile(f)
  else
    tex.error("LuaCJK: cannot find luacjk.lua")
  end
}
\group_end:
\directlua{
  luacjk.attr_font = \int_use:N \c_luacjk_font_attr_int
  luacjk.attr_cjk_space = \int_use:N \c_luacjk_cjk_space_attr_int
  luacjk.attr_xecglue = \int_use:N \c_luacjk_xecglue_attr_int
}
\AtBeginDocument
  {
    \luacjk_update_font:
    \g_luacjk_font_attr = \g_luacjk_fontid_int
  }
%    \end{macrocode}
%
% \begin{macrocode}
%</package>
%    \end{macrocode}
%
% \begin{macrocode}
%<*lua>
-- luacjk.lua
-- Copyright (C) 2026 by Mingyu Xia <myhsia@outlook.com>
-- This work may be distributed and/or modified under the conditions of the
-- LaTeX Project Public License, version 1.3c or later.
-- LuaCJK node processing.

luacjk = luacjk or {}

local M = luacjk
local GLYPH = node.id("glyph")
local GLUE = node.id("glue")
local HLIST = node.id("hlist")
local VLIST = node.id("vlist")
local new_node = node.new
local free_node = node.free
local has_attribute = node.has_attribute

M.attr_font = M.attr_font or nil
M.attr_cjk_space = M.attr_cjk_space or nil
M.attr_xecglue = M.attr_xecglue or nil

local function is_cjk (c)
  return (c >= 0x2E80 and c <= 0x9FFF)
      or (c >= 0xF900 and c <= 0xFAFF)
      or (c >= 0xFE30 and c <= 0xFE4F)
      or (c >= 0xFF00 and c <= 0xFFEF)
      or (c >= 0x20000 and c <= 0x3FFFF)
      or (c >= 0x3040 and c <= 0x30FF)
      or (c >= 0xAC00 and c <= 0xD7AF)
      or (c >= 0x1100 and c <= 0x11FF)
      or c == 0x2014 or c == 0x2018 or c == 0x2019
      or c == 0x201C or c == 0x201D or c == 0x2026
end

M.is_cjk = is_cjk

local function copy_glyph (n, fid)
  local g = new_node("glyph")
  g.font = fid
  g.char = n.char
  g.attr = n.attr
  g.lang = n.lang
  g.left = n.left
  g.right = n.right
  g.uchyph = n.uchyph
  g.xoffset = n.xoffset
  g.yoffset = n.yoffset
  g.expansion_factor = n.expansion_factor
  g.data = n.data
  g.subtype = n.subtype
  return g
end

local function rewrite_fonts (head)
  local prev = nil
  local cur = head
  while cur do
    local nxt = cur.next
    if cur.id == GLYPH and is_cjk(cur.char) and M.attr_font then
      local target = has_attribute(cur, M.attr_font)
      if target and target > 0 and target ~= cur.font then
        -- luaotfload inserts a private-use placeholder for a missing
        -- character.  Once the CJK glyph is replaced, that placeholder is
        -- no longer needed and would contribute a spurious advance width.
        if nxt and nxt.id == GLYPH and nxt.char == 0xF0000 then
          local after = nxt.next
          free_node(nxt)
          nxt = after
        end
        local g = copy_glyph(cur, target)
        g.next = nxt
        if prev then
          prev.next = g
        else
          head = g
        end
        free_node(cur)
        cur = g
      end
    end
    prev = cur
    cur = nxt
  end
  return head
end

local function get_glue (name)
  local s = tex.skip[name]
  if not s then
    return nil
  end
  return s
end

local function make_glue (name)
  local s = get_glue(name)
  if not s then
    return nil
  end
  local g = new_node("glue")
  g.width = s.width
  g.stretch = s.stretch
  g.shrink = s.shrink
  g.stretch_order = 0
  g.shrink_order = 0
  return g
end

local function set_glue (g, name)
  local s = get_glue(name)
  if not s then
    return false
  end
  g.width = s.width
  g.stretch = s.stretch
  g.shrink = s.shrink
  return true
end

local function spacing (head)
  local prev = nil
  local prevprev = nil
  local cur = head
  while cur do
    local nxt = cur.next
    if cur.id == GLYPH then
      local cur_cjk = is_cjk(cur.char)
      if prev and prev.id == GLYPH then
        local prev_cjk = is_cjk(prev.char)
        local name = nil
        if prev_cjk and cur_cjk then
          name = "l_luacjk_ccglue_skip"
        elseif prev_cjk or cur_cjk then
          name = "l_luacjk_ecglue_skip"
        end
        if name then
          local g = make_glue(name)
          if g then
            g.next = cur
            prev.next = g
            prevprev = prev
            prev = g
          end
        end
      elseif prev and prev.id == GLUE and prevprev and prevprev.id == GLYPH then
        local prev_cjk = is_cjk(prevprev.char)
        if prev_cjk and cur_cjk then
          local cjk_space = M.attr_cjk_space
            and (has_attribute(prev, M.attr_cjk_space) or 0) ~= 0
          if not cjk_space then
            prevprev.next = cur
            free_node(prev)
            prev = prevprev
          end
        elseif prev_cjk or cur_cjk then
          local xecglue = M.attr_xecglue
            and (has_attribute(prev, M.attr_xecglue) or 0) ~= 0
          if xecglue then
            set_glue(prev, "l_luacjk_ecglue_skip")
          end
        end
      end
    end
    if prev ~= cur then
      prevprev = prev
      prev = cur
    else
      prevprev = prevprev
    end
    cur = nxt
  end
  return head
end

local function process (head)
  head = rewrite_fonts(head)
  head = spacing(head)
  return head
end

M.process = process

function M.pre_linebreak (head, groupcode)
  return process(head)
end

function M.hpack (head, groupcode, size, packtype, direction, attributelist)
  return process(head)
end

luatexbase.add_to_callback("pre_linebreak_filter", M.pre_linebreak, "luacjk.spacing")
luatexbase.add_to_callback("hpack_filter", M.hpack, "luacjk.spacing")
%</lua>
%    \end{macrocode}
%
% \end{implementation}
%
% \Finale
%
\endinput
